% Options for packages loaded elsewhere
\PassOptionsToPackage{unicode}{hyperref}
\PassOptionsToPackage{hyphens}{url}
\documentclass[
]{article}
\usepackage{xcolor}
\usepackage[margin=1in]{geometry}
\usepackage{amsmath,amssymb}
\setcounter{secnumdepth}{-\maxdimen} % remove section numbering
\usepackage{iftex}
\ifPDFTeX
  \usepackage[T1]{fontenc}
  \usepackage[utf8]{inputenc}
  \usepackage{textcomp} % provide euro and other symbols
\else % if luatex or xetex
  \usepackage{unicode-math} % this also loads fontspec
  \defaultfontfeatures{Scale=MatchLowercase}
  \defaultfontfeatures[\rmfamily]{Ligatures=TeX,Scale=1}
\fi
\usepackage{lmodern}
\ifPDFTeX\else
  % xetex/luatex font selection
\fi
% Use upquote if available, for straight quotes in verbatim environments
\IfFileExists{upquote.sty}{\usepackage{upquote}}{}
\IfFileExists{microtype.sty}{% use microtype if available
  \usepackage[]{microtype}
  \UseMicrotypeSet[protrusion]{basicmath} % disable protrusion for tt fonts
}{}
\makeatletter
\@ifundefined{KOMAClassName}{% if non-KOMA class
  \IfFileExists{parskip.sty}{%
    \usepackage{parskip}
  }{% else
    \setlength{\parindent}{0pt}
    \setlength{\parskip}{6pt plus 2pt minus 1pt}}
}{% if KOMA class
  \KOMAoptions{parskip=half}}
\makeatother
\usepackage{graphicx}
\makeatletter
\newsavebox\pandoc@box
\newcommand*\pandocbounded[1]{% scales image to fit in text height/width
  \sbox\pandoc@box{#1}%
  \Gscale@div\@tempa{\textheight}{\dimexpr\ht\pandoc@box+\dp\pandoc@box\relax}%
  \Gscale@div\@tempb{\linewidth}{\wd\pandoc@box}%
  \ifdim\@tempb\p@<\@tempa\p@\let\@tempa\@tempb\fi% select the smaller of both
  \ifdim\@tempa\p@<\p@\scalebox{\@tempa}{\usebox\pandoc@box}%
  \else\usebox{\pandoc@box}%
  \fi%
}
% Set default figure placement to htbp
\def\fps@figure{htbp}
\makeatother
\setlength{\emergencystretch}{3em} % prevent overfull lines
\providecommand{\tightlist}{%
  \setlength{\itemsep}{0pt}\setlength{\parskip}{0pt}}
\usepackage{bookmark}
\IfFileExists{xurl.sty}{\usepackage{xurl}}{} % add URL line breaks if available
\urlstyle{same}
\hypersetup{
  pdftitle={paper},
  hidelinks,
  pdfcreator={LaTeX via pandoc}}

\title{paper}
\author{}
\date{\vspace{-2.5em}2026-03-18}

\begin{document}
\maketitle

\section{=========================================================}\label{section}

\section{INTRODUCTION}\label{introduction}

\section{=========================================================}\label{section-1}

Online ratings of films play an important role in shaping consumer
viewing decisions and influencing recommendation systems on streaming
platforms. It is therefore relevant to examine which characteristics of
audiovisual content are associated with higher audience ratings. One
potentially relevant characteristic is runtime.

Research by Moon et al.~(2009) suggests that longer films tend to
receive higher ratings on average, possibly because they are perceived
by viewers as more valuable and narratively complex. At the same time,
prolonged exposure to audiovisual stimuli may, according to Attention
Restoration Theory, lead to cognitive fatigue and reduced viewer
attention (Baumgartner \& Kühne, 2024). This may influence how viewers
evaluate the quality of a film or television production.

The relationship between runtime and ratings is also likely to differ by
genre, as different genres create distinct expectations regarding
narrative pacing, story development, and ideal length.

The central research question is therefore formulated as follows:

\textbf{To what extent does runtime predict the average IMDb rating of
films, and to what extent is this relationship moderated by genre?}

This research question aligns well with the available IMDb datasets and
focuses on whether movie length is associated with perceived quality, as
reflected in average IMDb ratings. In addition, the number of votes may
affect the reliability and stability of ratings. Movies with only a
small number of votes may receive more extreme or biased evaluations.
Therefore, the number of votes is included as a control variable.

The findings of this study are relevant for both consumers interpreting
online ratings and platforms that rely on rating-based recommendation
algorithms.

\section{=========================================================}\label{section-2}

\section{DATA}\label{data}

\section{=========================================================}\label{section-3}

The data used in this study are obtained from the
\href{https://datasets.imdbws.com/}{official IMDb datasets},
specifically we used these two datasets: -
\href{https://datasets.imdbws.com/title.basics.tsv.gz}{title.basics.tsv.gz}
-
\href{https://datasets.imdbws.com/title.ratings.tsv.gz}{title.ratings.tsv.gz}

These datasets are programmatically downloaded within the R script to
ensure full reproducibility. The title.basics dataset contains
\textbf{12,373,014} observations and includes information on title
characteristics such as title type, runtime, and genre. The
title.ratings dataset contains \textbf{1,650,357} observations and
provides the average user rating and the number of user votes. The two
datasets are merged using the unique identifier tconst.

To ensure a consistent and comparable sample, the analysis is restricted
to movies (titleType = ``movie''). Titles with missing values for
runtime, average rating, or number of votes are excluded.

Because IMDb assigns multiple genres to some titles, only the first
listed genre is retained as the primary genre classification to avoid
duplicate observations and ensure one observation per film. After
cleaning and filtering, the final merged dataset contains
\textbf{299,473} unique movie observations.

\section{=========================================================}\label{section-4}

\section{VARIABLES}\label{variables}

\section{=========================================================}\label{section-5}

\subsection{Dependent Variable: Average
Rating}\label{dependent-variable-average-rating}

The dependent variable is the average IMDb user rating of a movie,
measured by \texttt{averageRating}. It is recorded on a continuous scale
from 1 to 10 and reflects audience evaluations of the content. Higher
values indicate more favourable audience evaluations, while lower values
reflect more negative evaluations.

\subsection{Independent Variable:
Runtime}\label{independent-variable-runtime}

The independent variable is the runtime of a movie, measured by
\texttt{runtimeMinutes}. Runtime is expressed in minutes and treated as
a continuous variable. It represents the total length of the content. To
ensure a consistent and comparable sample, extreme runtime values are
excluded from the analysis. Movies shorter than 40 minutes and longer
than 240 minutes are removed, as these are unlikely to represent
standard feature films.

\subsection{Control Variable: Number of
Votes}\label{control-variable-number-of-votes}

The control variable is the total number of IMDb user votes, measured by
\texttt{numVotes}. This variable captures the popularity and visibility
of a title. Because the distribution of vote counts is strongly skewed,
the variable may be log-transformed in the regression analysis. This
log-transformation results in a more normally distributed variable
suitable for linear regression.

\subsection{Control Variable: Release
Year}\label{control-variable-release-year}

Release year of the movie is included as an additional control variable,
measured by \texttt{startYear}. This variable captures the year in which
the movie was released. The control variable is treated as a continuous
variable and it accounts for potential temporal trends in iMDb ratings,
as it may be that audience evaluation standards and rating behaviour may
have changed over time.

\subsection{Moderator: Genre}\label{moderator-genre}

Genre is included as a moderating variable to test whether the
relationship between runtime and rating differs across movie types.
Since IMDb allows multiple genres per title, only the first listed genre
is used as the primary genre classification in order to preserve one
observation per film.

Genre is treated as a categorical variable and operationalised using
dummy variables in the regression analysis, with one genre serving as
the reference category. For interpretability, genres are grouped into
the ten most frequent categories, with all remaining genres combined
into an \texttt{"Other"} category.

\section{=========================================================}\label{section-6}

\section{METHOD}\label{method}

\section{=========================================================}\label{section-7}

To answer the research question, linear regression analysis is employed,
as both the dependent variable (average IMDb rating) and the independent
variable (runtime in minutes) are continuous. Prior to analysis, the
data were prepared through several cleaning steps. The two IMDb datasets
(title.basics and title.ratings) were merged using the unique identifier
tconst and filtered to retain only feature films (titleType =
``movie''). Runtime outliers were removed by restricting observations to
films between 40 and 240 minutes, based on inspection of the empirical
quantile distribution. Titles with missing values on runtime, average
rating, number of votes, or release year were excluded via listwise
deletion, yielding a final analytic sample of 299,473 films.

To address the skewed distribution of the number of votes, this variable
was log-transformed (log\_votes = log(numVotes + 1)) before inclusion as
a control variable, improving the normality of its distribution and
reducing the influence of highly voted outliers. Because IMDb assigns
multiple genres per title, only the first-listed genre was retained as
the primary genre classification (main\_genre) to maintain one
observation per film. Genres were subsequently grouped into the ten most
frequent categories, with all remaining genres collapsed into an
``Other'' category, resulting in an eleven-level categorical moderator
(genre10).

The moderation hypothesis is tested by including interaction terms
between runtime and genre in the regression model, allowing the slope of
runtime to vary across genre categories. Release year was additionally
examined as a potential covariate through correlation analysis. All
analyses were conducted in R (version 4.0 or higher).

\section{=========================================================}\label{section-8}

\section{PREVIEW OF FINDINGS}\label{preview-of-findings}

\section{=========================================================}\label{section-9}

Preliminary analysis shows that the relationship between runtime and
average IMDb rating is statistically very weak. The Pearson correlation
between runtime and average rating is approximately r = 0.008, which is
very close to zero and indicates that longer movies do not
systematically receive higher ratings. This challenges the assumption
that longer films tend to be rated more favourably (Moon et al., 2009),
suggesting that runtime is not a reliable determinant of perceived
quality.

This pattern is also visible in the scatterplot below, where the
regression line shows only a slight positive trend and the data points
are widely dispersed. This suggests substantial unexplained variation in
ratings across movies with similar runtimes. In other words, runtime
alone does not appear to be a strong predictor of the average IMDb
rating of a movie.

\begin{figure}
\centering
\pandocbounded{\includegraphics[keepaspectratio]{../gen/output/runtime_vs_rating.png}}
\caption{Runtime vs rating}
\end{figure}

In other words, runtime alone does not appear to be a strong predictor
of the average IMDb rating of a movie.

The correlation between the log-transformed number of votes and average
rating is approximately \textbf{-0.08}, indicating a very small negative
relationship. This may reflect a diversification effect, whereby wider
viewership introduces more heterogeneous and critical evaluations.
Including the number of votes as a control variable is therefore
analytically important, as it isolates the effect of runtime from
variation in title visibility, thereby strengthening the internal
validity of the estimates.

As shown in the boxplot below. Average ratings vary substantially across
genres. For example, the average ratings of the genres \emph{Biography}
and \emph{Documentary} are noticeably higher than those of
\emph{Horror}. This suggests that genre may play an important moderating
role in the relationship between runtime and rating.

\begin{figure}
\centering
\pandocbounded{\includegraphics[keepaspectratio]{../gen/output/rating_by_genre.png}}
\caption{Rating by genre}
\end{figure}

Further analysis shows trends over time:

\begin{figure}
\centering
\pandocbounded{\includegraphics[keepaspectratio]{../gen/output/year_vs_rating.png}}
\caption{Year vs rating}
\end{figure}

Finally, we assess the validity of the regression model:

\begin{figure}
\centering
\pandocbounded{\includegraphics[keepaspectratio]{../gen/output/model_diagnostics.png}}
\caption{Model diagnostics}
\end{figure}

\section{=========================================================}\label{section-10}

\section{REPRODUCIBILITY}\label{reproducibility}

\section{=========================================================}\label{section-11}

The findings of this study are presented in a structured R Markdown
report. In this report, the complete data analysis is documented step by
step. The report includes data preparation, descriptive statistics,
regression analysis, and visualisations such as scatterplots and
boxplots.

Furthermore, the workflow is fully reproducible. The IMDb data can be
downloaded automatically and analysed through the script. This allows
other researchers to replicate the study easily or extend the analysis
using different selection criteria, such as focusing on a specific
genre. All scripts are version-controlled using Git and hosted in a
public GitHub repository, ensuring full transparency and accessibility
of the analytical workflow.

\section{=========================================================}\label{section-12}

\section{INTERPRETATION AND
RELEVANCE}\label{interpretation-and-relevance}

\section{=========================================================}\label{section-13}

The findings suggest that runtime, on its own, is a very weak predictor
of the average IMDb rating. This is relevant because it challenges the
assumption that longer movies necessarily receive higher ratings (Moon
et al., 2009). In practical terms, this indicates that longer runtime is
not a reliable strategy for achieving higher ratings.

In addition, the results show substantial variation in ratings across
genre categories. This implies that general statements about an ``ideal
movie length'' may be misleading and that such conclusions should
instead be approached in a genre-specific manner.

Including the number of votes as a control variable is essential for a
more accurate interpretation of the results. By controlling for the
number of votes, the effect of runtime is estimated while accounting for
variation in title visibility and popularity, which strengthens the
internal validity of the analysis. Including the release year as an
additional control variable further strenghtens the analysis by
accounting for potential changes in rating behaviour over time.
Together, these two control variables help ensure that the estimated
effect of runtime on average rating is not confounded by differences in
movie popularity or the time period in which a film was released.

\section{=========================================================}\label{section-14}

\section{REPOSITORY OVERVIEW}\label{repository-overview}

\section{=========================================================}\label{section-15}

The repository is organised to ensure reproducibility and transparency.
Scripts for data acquisition, preparation, and analysis are separated
into dedicated subdirectories, and the master makefile automates the
complete pipeline from data download to final output.

imdb-runtime-rating/ │ ├── data/ \# Raw and downloaded data files │ ├──
title.basics.tsv.gz \# IMDb basics dataset (auto-downloaded) │ └──
title.ratings.tsv.gz \# IMDb ratings dataset (auto-downloaded) │ ├──
reporting/ \# R Markdown reports │ ├── data\_exploration.Rmd \#
Exploratory data analysis and visualisations │ └── final\_report.Rmd \#
Full documented analysis report │ ├── src/ \# Source code │ ├──
analysis/ \# Statistical modelling scripts │ │ ├── analysis.R \# Linear
regression and moderation analysis │ │ └── makefile \# Pipeline
automation for analysis step │ │ │ └── data-preparation/ \# Data
cleaning and preparation scripts │ ├── data-cleaning.R \# Filtering,
merging, variable construction │ ├── download-data.R \# Automated
download of IMDb datasets │ └── makefile \# Pipeline automation for
preparation step │ ├── .gitignore \# Files excluded from version control
├── README.md \# Project documentation and running instructions └──
makefile \# Master pipeline automation

\section{=========================================================}\label{section-16}

\section{DEPENDENCIES}\label{dependencies}

\section{=========================================================}\label{section-17}

This project was developed using R (version 4.2.0 or higher) and
RStudio. The following R packages are required to run the workflow:

\begin{itemize}
\tightlist
\item
  \texttt{readr} for importing the IMDb .tsv.gz datasets
\item
  \texttt{dplyr} for data filtering, merging, and variable construction
\item
  \texttt{tidyr} for data reshaping and handling missing values
\item
  \texttt{rmarkdown} for rendering the R Markdown reports
\item
  \texttt{knitr} for document knitting and output generation
\end{itemize}

If any of these packages are not yet installed, they can be installed
collectively using:

\begin{verbatim}
install.packages(c("readr", "dplyr", "tidyr", "rmarkdown", "knitr"))
\end{verbatim}

\section{=========================================================}\label{section-18}

\section{RUNNING INSTRUCTIONS}\label{running-instructions}

\section{=========================================================}\label{section-19}

To reproduce the results of this workflow, follow these steps:

\textbf{Step 1: Clone the repository}

Using Command Prompt (Windows) or Terminal (Mac), clone the repository
to your local machine:

\begin{verbatim}
git clone https://github.com/course-dprep/The-effect-of-movie-runtime-on-average-IMDb-ratings.git
\end{verbatim}

Once cloned, all necessary folders and scripts will be available
locally.

\textbf{Step 2: Download the IMDb datasets}

Navigate to src/data-preparation/ and run download-data.R. This script
automatically creates a data/ folder and downloads the required IMDb
datasets:

\begin{verbatim}
# Create data folder
dir.create("data", showWarnings = FALSE, recursive = TRUE)

# Download IMDb datasets
download.file("https://datasets.imdbws.com/title.basics.tsv.gz",
              "data/title.basics.tsv.gz")

download.file("https://datasets.imdbws.com/title.ratings.tsv.gz",
              "data/title.ratings.tsv.gz")
\end{verbatim}

\textbf{Step 3: Load packages and read data}

Run data-cleaning.R to load the required packages and import the
downloaded datasets into RStudio:

\begin{verbatim}
# Load required libraries
library(readr)
library(dplyr)
library(tidyr)

# Read datasets
basics  <- read_tsv("data/title.basics.tsv.gz", na = "\\N")
ratings <- read_tsv("data/title.ratings.tsv.gz", na = "\\N")
\end{verbatim}

After completing these steps, the full analysis pipeline, including data
merging, cleaning, variable construction, and regression analysis can be
reproduced by knitting final\_report.Rmd in RStudio or by executing the
master makefile from the terminal. No manual data preprocessing is
required.

\section{=========================================================}\label{section-20}

\section{ABOUT}\label{about}

\section{=========================================================}\label{section-21}

This project is set up as part of the Master's course
\href{https://dprep.hannesdatta.com/}{Data Preparation \& Workflow
Management} at the
\href{https://www.tilburguniversity.edu/about/schools/economics-and-management/organization/departments/marketing}{Department
of Marketing}, \href{https://www.tilburguniversity.edu/}{Tilburg
University}, the Netherlands.

The project is implemented by team 1 members: - Bente van Brussel:
\href{mailto:b.j.m.vanbrussel@tilburguniversity.edu}{\nolinkurl{b.j.m.vanbrussel@tilburguniversity.edu}}
- Niels Deenen:
\href{mailto:n.deenen@tilburguniversity.edu}{\nolinkurl{n.deenen@tilburguniversity.edu}}
- David Lindwer:
\href{mailto:d.j.lindwer@tilburguniversity.edu}{\nolinkurl{d.j.lindwer@tilburguniversity.edu}}
- Demi Verburg:
\href{mailto:d.verburg@tilburguniveristy.edu}{\nolinkurl{d.verburg@tilburguniveristy.edu}}
- Marijn van Dooren:
\href{mailto:m.vandooren@tilburguniversity.edu}{\nolinkurl{m.vandooren@tilburguniversity.edu}}

\end{document}
